\section{Coefficient two and inert-prime rows on the same window}
\label{sec:coefficient-two-inert-window}

\paragraph{The full previous union is retained.}
Keep the actual host, original completions, prescribed pair,
page and finite datum of AC-G179-window.  Thus
\[
 \begin{gathered}
 H_u=\mathbb Q(\zeta_{2^{h(u)+2}},\zeta_7,
       \sqrt{-q},\sqrt{-d},\omega,\sqrt[3]{269}),\qquad
 K=\mathbb Q(\omega),\\
 u^{2+8\varepsilon}\leq2^{h(u)}\leq2u^{2+8\varepsilon},\quad
 0<\varepsilon<1/256,\quad T=\lfloor u^{399}\rfloor.
 \end{gathered}\tag{I2W.1}
\]
Write $Z^{\rm ref3}$ and $I_{399}^{\rm ref3}$ for the exact
omitted set and remainder of (R3W.5)--(R3W.8).
Retain $N_T,C(h),E_3(h),P(h),\Pi(h),C_\ell(h),B_T,A_T,R_T,S_T$
with precisely the meanings used there; in particular $N_T=T-1$.
The original residue embedding in the definition of $E_3(h)$ is
not changed by the new coefficient-two argument.

\paragraph{Two disjoint new sets.}
Let $\mathfrak D_h$ be all characters modulo $2^{h+2}$,
including the trivial one, with their actual primitive conductors.
Let
\[
 \Xi_7^*=\{1\}\sqcup
       \{\xi\bmod7:\operatorname{ord}\xi\in\{3,6\}\}.
 \tag{I2W.2}
\]
In particular $|\Xi_7^*|=5$; the quadratic character modulo
seven is not included.  Define
\[
 \begin{aligned}
 Z^2(u)=\{\,&(2,(\delta\xi)|_{G_K},e):
        \delta\in\mathfrak D_{h(u)},\ \xi\in\Xi_7^*,\ 1\leq e<T\},\\
 \mathfrak P_-(T)&=\{\ell\leq T:\ell\geq5\text{ prime},\
                      \ell\equiv2\pmod3,\ \ell\ne269\},\\
 Z^-(u)=\{\,&(\ell,1,e):\ell\in\mathfrak P_-(T),\
                   1\leq e<T,\ (\ell-1)\mid e\},\\
 J_-(T)&=\sum_{\ell\in\mathfrak P_-(T)}
                             \left\lfloor\frac{T-1}{\ell-1}\right\rfloor,\\
 Z^{\rm two,inert}(u)&=Z^{\rm ref3}(u)\sqcup Z^2(u)\sqcup Z^-(u).
 \end{aligned}\tag{I2W.3}
\]
All new finite labels have trivial Kummer cubic factor from
$K(\sqrt[3]{269})/K$ and trivial auxiliary $q,d$ factors.
The inert-prime family has the trivial finite label everywhere.
Let $I_{399}^{\rm two,inert}$ omit exactly the union (I2W.3)
from the original sum (TC.1).

\begin{theorem}[AC-G183-window: coefficient-two and inert-prime union]
\label{AC-G183-window}
Every new triple in (I2W.3) is an actual original-coefficient strict
zero.  The displayed unions are disjoint.  For $h\geq2$ and every
positive integer $T$,
\[
 \begin{aligned}
 |Z^2|&=5\cdot2^{h+1}N_T,& |Z^-|&=J_-(T),\\
 |Z^{\rm two,inert}\setminus Z^{\rm ref3}|
       &=5\cdot2^{h+1}N_T+J_-(T),\\
 |Z^{\rm two,inert}_3|&=(1+E_3(h))N_T,\\
 |Z^{\rm two,inert}_{269}|&=(8+C(h))N_T-B_T+12A_T-R_T-S_T.
 \end{aligned}\tag{I2W.4}
\]
For $\ell\geq5$, $\ell\ne269$, the coefficient-$\ell$ count is
\[
 \mathbf1_{\ell\in\Pi(h)}C_\ell(h)N_T+
 \mathbf1_{\ell\in\mathfrak P_-(T)}
                        \left\lfloor\frac{T-1}{\ell-1}\right\rfloor.
 \tag{I2W.5}
\]
Consequently the cumulative count is
\[
 \begin{split}
 |Z^{\rm two,inert}|={}&
  (9+C(h)+E_3(h)+P(h)+5\cdot2^{h+1})N_T\\
 &-B_T+12A_T-R_T-S_T+J_-(T).
 \end{split}\tag{I2W.6}
\]
The inclusion of omitted sets is proper when $T\geq2$.
This does not assert strict decrease of the numerical remainder.

For the normalized positive-index cost $Q_u$ of this entire union,
\[
 \begin{aligned}
 I_{399}&=I_{399}^{\rm two,inert}+Q_u,\\
 I_{399}^{\rm two,inert}&\leq I_{399}^{\rm ref3}\leq I_{399},\\
 0\leq Q_u&\leq8D_u^2T^2\log(8TD_u)
          +\tfrac12D_u\mathcal R_uT^2+2D_ur_uc_0T\\
 &=O_\varepsilon(u^{802+16\varepsilon}(\log u)^2).
 \end{aligned}\tag{I2W.7}
\]
One complete (FL.4) majorant covers the whole union, including
coefficient two and the growing set of inert coefficient primes.
There is no additional factor for their numbers or character counts.
The identities use nonnegative extended sums and do not assert
that any individual positive derived index is zero.
The positive-relative-slack premise for this remainder supplies
the full original premise on the identical common window, with
all original coefficients, full local conditions and actual $R.t$.
\end{theorem}

\begin{proof}
\proofstep{1. Actual coefficients and the two different strict arguments.}
The host contains both cyclotomic factors appearing in (I2W.2).
Every $\delta\in\mathfrak D_h$ has order dividing $2^h$;
every $\xi\in\Xi_7^*$ has order dividing six.
Their products therefore have order dividing $3\cdot2^h$, so
their values already lie in the prescribed $O_2$ of (FC.2).
For every coefficient prime in use the identity is
\[
 (O_\ell\otimes_{\mathbb Z_\ell}\mathbb Q_\ell/\mathbb Z_\ell)
               (\eta)(-e)
       =(\operatorname{Frac}(O_\ell)/O_\ell)(\eta c_\ell^{-e}).
 \tag{I2W.8}
\]
This is the original lattice and completion, including its
ramification, not a change of coefficient field.

For coefficient two and $\xi=1$, use (G182.9) with $\kappa=1$
on this original $O_2$ for every pure two-power $\delta$, including
the trivial label, and every positive Tate weight.
That theorem uses the full coefficient's pro-two action and the
unique prime of $K$ above two; it is not a residual averaging
argument by the order of the growing host.
For $\operatorname{ord}\xi=3$ or six, use (G184.1)--(G184.2)
with $s=0$, $\kappa=1$, namely its literal strict theorem with the same $O_2$, $\delta,\xi$ and weight.
Its class-field restriction is applied only to the residual
order-three character.  In particular, the full coefficient with
an order-six $\xi$ is not asserted unramified away two after
restriction to $K\mathbb Q(\zeta_7)^+$.
The nontrivial order-three residue on inertia at seven and at
the local Frobenius over two supplies the original-uniformizer
injections on $K$ in that theorem.

For $\ell\in\mathfrak P_-(T)$ and $(\ell-1)\mid e$, use
(G182.10) with $\kappa=1$ at the actual original
$O_\ell$, with trivial finite label.
The prime is inert in $K$ and hence unique over $\ell$.
The divisibility of the weight makes the finite coefficient
quotients' cyclotomic action pro-$\ell$; no finite character at
the coefficient prime is added.

In every case the required strict conditions are full local zero
above the coefficient prime and zero in away inertia.
Inflation from the prescribed ramification group in (FC.2) to
$G_K$ is injective on the same coefficient module.
An inflated class is zero on all outside-set inertia, and inside
the set it satisfies precisely the stated full local conditions.
It therefore injects into the literal zero group from the new
strict theorems.  All old zeros remain supplied by AC-G179-window.

\proofstep{2. Exact sets and finite counts, including all overlaps.}
The character group modulo $2^{h+2}$ has $2^{h+1}$ elements,
and the character group modulo seven is cyclic of order six.
It has two characters of order three and two of order six.
The coprime-conductor product distinguishes every pair
$(\delta,\xi)$ in (I2W.3).
Restriction to $G_K$ preserves this distinctness: a nontrivial
ratio supported only at two and seven cannot equal $\chi_{-3}$.
Thus there are exactly $5\cdot2^{h+1}$ new coefficient-two labels.
The old union has no coefficient-two triples, so no overlap occurs.

At an inert prime the admissible weights are the positive
multiples of $\ell-1$ below $T$; their number is
$\lfloor(T-1)/(\ell-1)\rfloor$.
There are none unless $\ell\leq T$, making $J_-(T)$ finite.
This bound is a consequence of the positive-weight range, not
a new cutoff on the original all-prime sum.
At $269$ the old trivial label already covers every weight,
so that prime is explicitly excluded from $\mathfrak P_-(T)$.
At every other $\ell\geq5$ the old G176/G179 labels have a
nontrivial primitive high two-power factor, whose conductor
cannot cancel a factor modulo seven.  None is the trivial label.
Consequently $Z^-$ is disjoint from the entire old union and
from $Z^2$.  Its count is exactly $J_-(T)$, including primes
that were not in the previous finite set $\Pi(h)$.
Adding the unchanged (R3W.6) counts proves (I2W.4)--(I2W.6).

\proofstep{3. One full local error, also at coefficient two.}
Keep the original Tate-dual codifferent and ranks $b_\ell$.
The determinant inequality (EU.2), expressly applied at all
coefficient primes by AC-K26-elliptic, gives for each omitted triple
\[
 \frac{\log\#k_{O_\ell}}{b_\ell}\max(0,i_{\ell,\eta,e})
 \leq\frac{|\log\#H^2(C_{\ell,\eta,e})
                       -\log\#\mathcal S_{\ell,\eta}(e)|}{b_\ell}.
 \tag{I2W.9}
\]
No integral Fourier projector or division by $|G_u|$ is used.
The existing (FC.4)--(FC.5) comparison costs at two and three
remain in the retained interface.
At each weight the entire new union is one nonnegative subsum
of (FL.4), whose statement includes coefficient two.
Sum once over $1\leq e<T$ and use $\sum e\leq T^2/2$ to get
(I2W.7).  The unchanged bounds
$D_u\leq288u^{2+8\varepsilon}$, $r_u\leq12$ and
$\mathcal R_u=O_\varepsilon(u^{\varepsilon/5}(\log u)^2)$
give the displayed exponent uniformly in the prescribed pair.
Nonnegative additivity gives the extended-sum decomposition
without subtracting infinities.

\proofstep{4. The identical common window and remaining premise.}
Fix the retained finite datum, page and $\varepsilon$.
Assume the same (TC.3) quantifiers for the new remainder:
$J\geq1$, $X_0$ and $s\in(0,1)$ are independent of
$X,q,d,u,\ell,\eta,e$, and every prescribed pair at $X\geq X_0$
has one $X^{1/J}\leq W$, $W^2\leq X$, on which
$I_{399}^{\rm two,inert}(H_u,u)\leq(1-s)\exp(u^\varepsilon/2)$
throughout $[W,W^2]$.
The single uniform error in (I2W.7) is
$o(\exp(u^\varepsilon/2))$.
Choose $U_0\geq2$ making it at most half the slack, raise
$X_0$ to $\max(X_0,U_0^J)$, and keep exactly the supplied $W$.
The full original premise holds there with slack $s/2$.
Every older premise implies this new one by nonnegativity.
Uniform little-$o$ and fixed finite simultaneous families use
the same argument and a maximum of the uniform constants;
no independent windows are intersected.
All original coefficients, full local conditions, actual $R.t$
and simultaneous growing-host/prime/character/weight quantifiers
are retained.  No stronger AW.1 hypothesis is introduced.
\end{proof}

\paragraph{Boundary and retained dictionary.}
The quadratic modulo-seven factor by itself is not added at
coefficient two; neither are nontrivial cubic or auxiliary
$q,d$ factors.  Unsupported prime/weight/character combinations
remain in the remainder.
The numerical remainder, mixed theta/Selmer realization
$\mathcal R_{\rm cub}$ and H04.4 remain open.
Only the already audited SGA~1, Expos\'e~V, \S\S4,7 and
Stacks~03PK finite-cover/Artin--Tate dictionary is retained.
No new motive, source archive scan or Lean execution is claimed.
